\section{Conclusion}
\label{sec:conclusion}

We have shown that, in a cross-lingual FLORES+ setup with an XLM-RoBERTa-base backbone, token-level anisotropy is a measurable property of vanilla multilingual ColBERT and one that can be reshaped with a single regularizer term added to InfoNCE training. The regularizer has one hyperparameter and zero inference cost; it improves token-embedding geometry monotonically with regularization strength and recovers $+3.45$ percentage points of MaxSim P@1 on English-Swahili and $+1.99$ percentage points on English-Arabic, with the gains extending zero-shot to two further low-resource languages (Nepali $+1.46$, Tamil $+0.88$). Mechanism ablations against active-token and cross-attention placements show that, of the three placements we tested, the uniform variant is the only one that produces both the geometric reshaping and the retrieval gain. We caution that all of these claims are tied to our specific evaluation setting; whether the same geometric pathology drives retrieval errors on harder benchmarks (MIRACL, mMARCO) or on larger and retrieval-specialized backbones (multilingual E5) remains an empirical question for future work. Within the setting we evaluate, however, token geometry is a tractable lever for low-resource multilingual retrieval.
